\chapter{Immunoglobulins (IgA, IgG, IgM)}

\section*{What Is This Test?}
Quantitative immunoglobulin testing measures the serum concentrations of the three major classes of immunoglobulins (antibodies): IgG, IgA, and IgM. Immunoglobulins are glycoprotein molecules produced by B cells and plasma cells that function as the effector arm of humoral immunity. IgG is the most abundant immunoglobulin in serum (75--80\% of total immunoglobulins) and provides long-term protection against pathogens through four subclasses (IgG1, IgG2, IgG3, IgG4), each with distinct functional properties. IgG is the only immunoglobulin that crosses the placenta, providing passive immunity to the fetus and newborn. IgA is the major immunoglobulin in mucosal secretions (respiratory, gastrointestinal, and genitourinary tracts) and exists as two subclasses (IgA1, IgA2); serum IgA is predominantly monomeric, while mucosal IgA is dimeric with secretory component. IgM is the first antibody produced in an immune response (primary response), is a potent activator of the classical complement pathway, and exists as a pentamer in serum. Quantitative immunoglobulin measurement is essential for the diagnosis and classification of immunodeficiencies (primary and secondary), monitoring of plasma cell dyscrasias (multiple myeloma, Waldenstrom macroglobulinemia), and assessment of humoral immune function.

\section*{Alternative Names}
\begin{itemize}
  \item Immunoglobulins
  \item Quantitative Immunoglobulins
  \item IgG, IgA, IgM
  \item Serum Immunoglobulins
  \item Ig Levels
  \item Total Immunoglobulins
  \item Gamma Globulins
  \item Quantitative Ig
\end{itemize}
\section*{Principle}
Serum IgG, IgA, and IgM concentrations are measured by immunoturbidimetry or immunonephelometry on automated clinical chemistry analyzers. Specific anti-human IgG, IgA, or IgM antibodies are incubated with patient serum. The formation of antigen-antibody complexes results in light scattering (turbidity) or light diffraction (nephelometry), which is measured photometrically or by nephelometry. The concentration of each immunoglobulin is proportional to the extent of light scattering and is determined from a standard curve generated using calibrators of known concentration traceable to international reference preparations (e.g., CRM 470/RPPHS, ERM-DA470k). Total protein level and serum protein electrophoresis (SPEP) provide complementary information by showing the distribution of proteins including the gamma globulin (immunoglobulin) fraction. IgG subclass quantitation (IgG1, IgG2, IgG3, IgG4) is performed by radial immunodiffusion or ELISA when IgG subclass deficiency is suspected. Free light chain assays are performed separately by immunonephelometry using antibodies specific for kappa and lambda free light chains.

\section*{History}
The discovery of immunoglobulins began with Arne Tiselius and Elvin Kabat in 1939, who demonstrated by electrophoresis that antibodies migrated in the gamma globulin fraction of serum. The distinct classes of immunoglobulins were identified through immunochemical techniques in the 1950s--1960s: IgM was discovered in 1939 (Waldenstrom macroglobulinemia), IgG was characterized in detail by Rodney Porter and Gerald Edelman (Nobel Prize, 1972), IgA was identified as a distinct class in 1953, IgD in 1965, and IgE in 1966--1967. IgG subclasses were described in the 1960s. The development of nephelometry and turbidimetry in the 1970s--1980s enabled rapid, automated measurement of immunoglobulins, replacing labor-intensive radial immunodiffusion. The primary immunodeficiency diseases were classified by the World Health Organization beginning in 1970, with updated classifications by the International Union of Immunological Societies (IUIS) every few years \cite{picard2018international}.

\section*{Why This Test Is Ordered}
\begin{itemize}
  \item Screening for primary immunodeficiency diseases in patients with recurrent sinopulmonary infections (particularly sinopulmonary infections with encapsulated organisms including S. pneumoniae and H. influenzae), unexplained bronchiectasis, or recurrent gastrointestinal infections \cite{bonilla2015practice}
  \item Diagnosis and classification of common variable immunodeficiency (CVID), the most common symptomatic primary immunodeficiency, characterized by low IgG and IgA with variable IgM
  \item Evaluation of selective IgA deficiency, the most common primary immunodeficiency (prevalence 1:500 in Caucasians), characterized by absent or near-absent IgA with normal IgG and IgM
  \item Diagnosis of hyper-IgM syndromes, X-linked agammaglobulinemia (Bruton disease), and severe combined immunodeficiency (SCID)
  \item Screening for plasma cell dyscrasias: multiple myeloma (IgG myeloma most common, followed by IgA, with suppressed normal immunoglobulins), Waldenstrom macroglobulinemia (IgM), and monoclonal gammopathy of undetermined significance (MGUS) \cite{rajkumar2006monoclonal}
  \item Monitoring response to immunoglobulin replacement therapy in patients with primary immunodeficiency
  \item Evaluation of polyclonal hypergammaglobulinemia in autoimmune diseases (SLE, Sjogren syndrome, sarcoidosis), chronic infections (HIV, chronic hepatitis, infective endocarditis), and chronic liver diseases (autoimmune hepatitis, primary biliary cholangitis, alcoholic cirrhosis)
  \item Assessment of humoral immune function in secondary immunodeficiency states including HIV/AIDS, chronic lymphocytic leukemia (CLL), lymphoma, multiple myeloma, nephrotic syndrome, and protein-losing enteropathy
  \item Pre-treatment evaluation of patients prior to B-cell-depleting therapy (rituximab) or during immunosuppressive treatment
\end{itemize}

\section*{Primary vs Secondary Test}
This is a primary diagnostic test for humoral immunodeficiency diseases and plasma cell dyscrasias. Quantitative immunoglobulins are the first-line laboratory investigation when a humoral immune defect is suspected. In the context of monoclonal gammopathy, quantitative immunoglobulins complement serum protein electrophoresis (SPEP), immunofixation electrophoresis (IFE), and serum free light chain assay.

\section*{How This Test Is Performed}
For most patients, a health care professional will take a blood sample from a vein in your arm using a small needle. After the needle is inserted, a small amount of blood will be collected into a test tube or vial. You may feel a little sting when the needle goes in or out. This usually takes less than five minutes.

\section*{What Preparation Is Needed}
No special preparation is required. Fasting is not necessary. Patients receiving immunoglobulin replacement therapy (IVIG or SCIG) should ideally have trough levels measured immediately before the next infusion. Immunoglobulin levels are affected by drugs including corticosteroids, immunosuppressants (cyclophosphamide, mycophenolate, azathioprine), chemotherapy, rituximab and other B-cell-depleting agents, antiepileptic medications (phenytoin, carbamazepine), disease-modifying antirheumatic drugs (DMARDs), and captopril.

\section*{Contraindications and Precautions}
\begin{itemize}
  \item No absolute contraindications. Specimen should be collected in a serum separator tube (SST, gold-top) or red-top tube. Grossly hemolyzed or lipemic specimens may interfere with nephelometric measurements. Immunoglobulin reference ranges are age-dependent --- newborns have low IgG levels at birth (receiving maternal IgG transplacentally); IgG reaches a physiologic nadir at 3--6 months of age (transient hypogammaglobulinemia of infancy) as maternal IgG is catabolized and the infant's own IgG production is still ramping up. IgA and IgM production develops more gradually, with adult levels reached by 6--16 years. Immunoglobulin levels during pregnancy are slightly lower due to hemodilution. Malnutrition, nephrotic syndrome, and protein-losing enteropathy cause secondary hypogammaglobulinemia through protein loss or decreased synthesis. Dextran and radiopaque contrast agents may cause falsely elevated nephelometric results. Monoclonal immunoglobulins may interfere with nephelometric measurement of polyclonal immunoglobulins --- SPEP and IFE should always be correlated.
\end{itemize}
\section*{Risks}
Minimal risk. Slight pain or bruising at the venipuncture site. Rare: hematoma, infection, vasovagal syncope.

\section*{What Happens After the Test}
Results are typically available within 1 to 2 hours. Low immunoglobulin levels in a patient with recurrent infections suggest primary or secondary immunodeficiency and may warrant further testing including IgG subclass measurement, specific antibody response to vaccines (tetanus, diphtheria, pneumococcal polysaccharide vaccine), flow cytometry for lymphocyte subsets, and referral to an immunologist. An elevated immunoglobulin level of a single class with suppressed other immunoglobulins raises concern for a monoclonal gammopathy and should prompt serum protein electrophoresis (SPEP), immunofixation electrophoresis (IFE), and serum free light chain assay. Polyclonal hypergammaglobulinemia (diffuse increase in gamma globulins with normal electrophoretic pattern) suggests chronic inflammation, autoimmune disease, or chronic liver disease.

\section*{Understanding Results}

\subsection*{Below Normal}
\begin{itemize}
  \item Common variable immunodeficiency (CVID): The most common symptomatic primary immunodeficiency presenting in the second or third decade of life. Characterized by low IgG AND low IgA (or low IgM) with impaired antibody responses to vaccination. Patients present with recurrent sinopulmonary infections, autoimmune cytopenias (ITP, AIHA), and an increased risk of lymphoma \cite{cunninghamrundles1999common}
  \item Selective IgA deficiency: IgA <7 mg/dL with normal IgG and IgM. The most common primary immunodeficiency (1:500 in Caucasians). Many individuals are asymptomatic; those with symptoms present with recurrent sinopulmonary infections, autoimmune diseases, and allergic disorders. Patients have an increased risk of anaphylaxis to blood products due to anti-IgA antibodies
  \item IgG subclass deficiency: Normal total IgG with deficiency of one or more IgG subclasses (IgG1 deficiency is most clinically significant due to impaired antibody response to protein antigens). Diagnosed by IgG subclass quantitation when total IgG is normal but clinical suspicion for immunodeficiency is high
  \item X-linked (Bruton) agammaglobulinemia: Profoundly low or absent all immunoglobulin classes due to a defect in the BTK gene blocking B-cell development. Presents in infancy after maternal IgG wanes (4--12 months) with recurrent pyogenic infections
  \item Severe combined immunodeficiency (SCID): Absent or profoundly low all immunoglobulin classes with absent T-cell function. Presents in early infancy with severe infections (Pneumocystis jirovecii pneumonia, chronic diarrhea, failure to thrive). Universal newborn screening for SCID uses T-cell receptor excision circles (TREC) quantification
  \item Transient hypogammaglobulinemia of infancy: Physiologic IgG nadir at 3--6 months with delayed IgG production. Spontaneous recovery typically by 2--4 years of age
  \item Secondary hypogammaglobulinemia: Nephrotic syndrome (urinary IgG loss), protein-losing enteropathy (GI IgG loss), severe burns, malnutrition, and chronic lymphocytic leukemia (CLL) with impaired B-cell function. Immunosuppressive medications (corticosteroids, cyclophosphamide, rituximab) lower immunoglobulin levels
  \item Multiple myeloma with suppressed normal immunoglobulins (immunoparesis): Monoclonal immunoglobulin elevation with suppression of uninvolved immunoglobulins (e.g., IgG kappa myeloma with low IgA and IgM) due to disruption of normal plasma cell function
\end{itemize}

\subsection*{Above Normal}
\begin{itemize}
  \item Polyclonal hypergammaglobulinemia (IgG >1,600 mg/dL): Chronic infections (HIV/AIDS, chronic hepatitis B and C, infective endocarditis, tuberculosis), autoimmune diseases (systemic lupus erythematosus, Sjogren syndrome, rheumatoid arthritis, sarcoidosis), and chronic liver diseases (autoimmune hepatitis, primary biliary cholangitis, alcoholic cirrhosis). Polyclonal increases show a broad gamma band on SPEP
  \item IgG multiple myeloma: Monoclonal IgG elevation (>3,000 mg/dL) with a sharp peak (M-spike) on SPEP. Suppressed IgA and IgM (immunoparesis). Patients present with CRAB features: hyperCalcemia, Renal insufficiency, Anemia, Bone lesions. IgG myeloma accounts for approximately 50--55\% of all multiple myelomas
  \item IgA multiple myeloma: Monoclonal IgA elevation with M-spike on SPEP. Accounts for approximately 20--25\% of multiple myelomas. IgA migrates in the beta or gamma region on SPEP
  \item IgM monoclonal gammopathy: Waldenstrom macroglobulinemia (IgM monoclonal gammopathy with lymphoplasmacytic lymphoma infiltration of bone marrow). Presents with hyperviscosity syndrome (headaches, blurred vision, epistaxis, neurologic symptoms) and constitutional B-symptoms. IgM MGUS (monoclonal gammopathy of undetermined significance) has a lower serum IgM level and no clinical symptoms
  \item Light chain multiple myeloma: Normal or low IgG, IgA, and IgM with elevated serum free light chains (kappa or lambda) and a high kappa/lambda ratio. Accounts for approximately 15--20\% of myelomas. SPEP may not show an M-spike; serum free light chain assay and 24-hour urine protein electrophoresis are critical
  \item Polyclonal IgA elevation: Chronic liver disease (alcoholic cirrhosis and chronic hepatitis produce IgA excess), IgA nephropathy (Berger disease), and infections of mucosal surfaces (respiratory and gastrointestinal)
  \item Polyclonal IgM elevation: Primary biliary cholangitis, acute viral infections (EBV, CMV, acute hepatitis A), and congenital infections (toxoplasmosis, rubella, CMV, herpes simplex)
  \item IgG4-related disease: Elevated serum IgG4 levels (>135 mg/dL) in the context of tumefactive lesions, organ enlargement, storiform fibrosis, obstructive phlebitis, and lymphoplasmacytic infiltration with IgG4-positive plasma cells
\end{itemize}

\section*{Normal Results}
\begin{itemize}
  \item \textbf{IgG:}
  \begin{itemize}
    \item Adults: 700--1,600 mg/dL
    \item Newborns: 700--1,400 mg/dL (predominantly maternal IgG transplacentally)
    \item Infants 3--6 months: 200--600 mg/dL (physiologic nadir)
    \item Children 1--5 years: 400--1,200 mg/dL
    \item Children 6--14 years: 500--1,500 mg/dL
  \end{itemize}
  \item \textbf{IgA:}
  \begin{itemize}
    \item Adults: 70--400 mg/dL
    \item Newborns: 0--5 mg/dL (IgA does not cross the placenta)
    \item Children 1--5 years: 10--100 mg/dL
    \item Children 6--14 years: 30--200 mg/dL
  \end{itemize}
  \item \textbf{IgM:}
  \begin{itemize}
    \item Adults: 40--230 mg/dL
    \item Newborns: 5--30 mg/dL (IgM does not cross the placenta; elevated in congenital infections)
    \item Children 1--5 years: 30--120 mg/dL
    \item Children 6--14 years: 40--150 mg/dL
  \end{itemize}
  \item \textbf{IgG subclasses (adults):} IgG1 280--1,020 mg/dL; IgG2 120--740 mg/dL; IgG3 20--170 mg/dL; IgG4 10--130 mg/dL
  \item \textbf{IgE:} <100 IU/mL (measured separately for allergy evaluation)
  \item Results must be interpreted with age-appropriate reference ranges and correlated with serum protein electrophoresis and immunofixation electrophoresis when monoclonal gammopathy is suspected
\end{itemize}

\section*{Follow-up Tests}
\begin{itemize}
  \item \textbf{Serum protein electrophoresis (SPEP):} Separates serum proteins into albumin, alpha-1, alpha-2, beta, and gamma fractions. Detects monoclonal proteins (M-spike) as a sharp band in the gamma (or beta) region. Quantifies the M-protein concentration. Essential for the evaluation of monoclonal gammopathies
  \item \textbf{Serum immunofixation electrophoresis (IFE) or immunotyping:} Identifies the immunoglobulin heavy chain class (IgG, IgA, IgM, IgD, IgE) and light chain type (kappa or lambda) of a monoclonal protein detected by SPEP. More sensitive than SPEP for detecting small monoclonal proteins
  \item \textbf{Serum free light chain (FLC) assay:} Measures kappa and lambda free light chains and calculates the kappa/lambda ratio. Essential for the diagnosis of light chain multiple myeloma, AL amyloidosis, and light chain deposition disease. The free light chain assay combined with SPEP and IFE is the recommended screening panel for monoclonal gammopathies
  \item \textbf{24-hour urine protein electrophoresis (UPEP) and urine immunofixation:} Detects Bence Jones protein (monoclonal free light chains excreted in the urine), a hallmark of multiple myeloma
  \item \textbf{IgG subclass measurement (IgG1, IgG2, IgG3, IgG4):} If total IgG is normal but clinical suspicion for immunodeficiency is high (recurrent sinopulmonary infections with normal total IgG)
  \item \textbf{Antibody response to immunization:} Tetanus and diphtheria toxoid IgG (protein antigens) and pneumococcal polysaccharide IgG (polysaccharide antigens) before and 4--6 weeks after vaccination to assess functional antibody production. Impaired antibody response confirms humoral immunodeficiency
  \item \textbf{Flow cytometry for lymphocyte subsets (B cells, T cells, NK cells):} Quantifies and characterizes lymphocyte populations. Absent B cells are characteristic of X-linked agammaglobulinemia; low T cells are characteristic of SCID
  \item \textbf{Bone marrow biopsy and aspiration:} If monoclonal gammopathy is suspected or diagnosed, to assess plasma cell percentage (>=10\% for symptomatic myeloma), morphology, and clonality by flow cytometry or immunohistochemistry \cite{rajkumar2014international}
  \item \textbf{Skeletal survey (or whole-body low-dose CT or PET-CT):} To detect lytic bone lesions in multiple myeloma
  \item \textbf{Serum calcium, creatinine, complete blood count:} To assess for CRAB criteria (hyperCalcemia, Renal insufficiency, Anemia, Bone lesions) that define symptomatic myeloma requiring treatment
\end{itemize}

\section*{References}
\bibliographystyle{unsrtnat}
\bibliography{references}

\clearpage
\section*{Regulatory Status and Authorization Date}
This chapter describes a test category, not a specific commercial product. FDA, CDSCO, EU IVDR, and MHRA approval or registration dates therefore cannot be assigned without the assay manufacturer, product name, instrument, and intended use. For laboratory-developed tests, an individual agency approval date may not exist. See the Preface for the interpretation of regulatory status.

\clearpage
